%% Manuscript: Site Effects in Multi-Centre EEG Data for Neuroscreening
%% Target journal: Biomedical Signal Processing and Control (Elsevier)
%% LaTeX class: elsarticle, numerical citations
%%
%% Compile with:
%%   pdflatex manuscript_eeg_harmonization.tex
%%   bibtex   manuscript_eeg_harmonization
%%   pdflatex manuscript_eeg_harmonization.tex
%%   pdflatex manuscript_eeg_harmonization.tex

\documentclass[preprint,12pt]{elsarticle}

%% For final journal submission use:
%% \documentclass[final,3p,times]{elsarticle}

\usepackage{amsmath}
\usepackage{amssymb}
\usepackage{graphicx}
\usepackage{booktabs}
\usepackage{multirow}
\usepackage{hyperref}
\usepackage{lineno}
\usepackage{microtype}
\usepackage{xcolor}
\usepackage{subcaption}

%% Uncomment for review mode:
% \linenumbers

\journal{Biomedical Signal Processing and Control}

% -----------------------------------------------------------------------
\begin{document}

\begin{frontmatter}

\title{Harmonisation of Multi-Site EEG Data for Clinical Neuroscreening:
       Quantifying the Site Effect and Comparing Correction Strategies}

\author[aff1]{Kacper Chorzela}
\author[aff1]{Julia Pelczar}
\author[aff1]{Jaros\l{}aw {\.Z}ygierewicz\fnref{fn1}\corref{cor1}}
\ead{jarekz@fuw.edu.pl}

\cortext[cor1]{Corresponding author.}
\fntext[fn1]{ORCID: 0000-0002-7536-0735}

\affiliation[aff1]{organization={Faculty of Physics, University of Warsaw},
             addressline={Pasteura 5},
             city={Warsaw},
             postcode={02-093},
             country={Poland}}

% -----------------------------------------------------------------------
\begin{abstract}
Pooling electroencephalographic (EEG) recordings from multiple clinical
centres introduces systematic, non-biological variance---commonly
termed the \emph{site effect} or \emph{batch effect}---that can
severely degrade the generalisation of machine-learning diagnostic
models.
We investigate the site effect in a large-scale collection of
54\,779 clinical EEG recordings from 30 Polish hospitals
(hereafter $\mathrm{ELM}_s$), and evaluate a range of harmonisation
strategies spanning simple site-wise standardisation, the neuroCombat
and CovBat algorithms, and PCA-based feature-space
pre-orthogonalisation applied as a preprocessing step before any
marginal correction method.
Site-effect characterisation and harmonisation evaluation are
conducted on the subset of normal recordings within $\mathrm{ELM}_s$
(hereafter $\mathrm{ELM}_n$), while downstream neuroscreening
performance is assessed on the full $\mathrm{ELM}_s$ corpus.

A two-classifier evaluation framework is employed throughout: a
logistic regression linear probe and a CatBoost gradient-boosted
nonlinear classifier are applied symmetrically to every condition,
including both hand-crafted feature vectors and the recording-level
representations produced by adversarial neural models.
This framework separates the linearly accessible component of the
site effect from the nonlinear covariance structure that survives
marginal correction.

The nonlinear classifier trained on the unharmonised
2\,850-dimensional feature set (Riemannian covariance projections,
multitaper power spectra, and inter-electrode coherence estimates)
achieved an overall Matthews Correlation Coefficient (MCC) of
$0.865 \pm 0.003$, confirming that site identity is strongly encoded
in the feature space.
SHAP explainability analysis identified inter-electrode coherence
features---in particular occipital O1--O2 coherence and beta-band
temporal coherence---as the primary carriers of site-discriminative
information, implicating differences in recording protocols and
patient-preparation procedures across institutions.

Applying marginal correction methods without pre-orthogonalisation
consistently \emph{increased} the nonlinear site-classification MCC
towards a near-ceiling value, demonstrating a \emph{harmonisation
paradox}: feature-by-feature correction reorganises rather than
eliminates the multivariate site structure encoded in inter-feature
covariances.
The linear probe, by contrast, assigned near-chance MCC to all
harmonised conditions, establishing that marginal correction
successfully removes the \emph{linear} component of the site
effect whilst leaving its nonlinear covariance component intact.
PCA pre-orthogonalisation substantially attenuates the paradox and,
in combination with any of the marginal correction methods, yields
more uniform cross-site neuroscreening performance than marginal
correction alone.

Signal-level adversarial harmonisation via a Domain Adversarial
Neural Network (DANN) integrated with the multiple-instance learning
(MIL) architecture of Poziomska~et~al.~\cite{poziomska2026quantity}
is evaluated as a complementary direction.
Four DANN variants---MINet (signal-level baseline without domain
adversarial training), MINet-MTL (multi-task baseline without gradient
reversal), MINet-DANN (full adversarial training), and
MINet-DANN+tuning (adversarial training with post-hoc classifier
recalibration on a frozen backbone)---are assessed under the same
symmetric two-classifier framework and the same Leave-One-Site-Out
(LOSO) protocol.

\textbf{[PLACEHOLDER: Complete aggregate results for all conditions
once all DANN experiments are finalised.]}
\end{abstract}

\begin{keyword}
EEG \sep site effect \sep batch effect \sep data harmonisation \sep
ComBat \sep CovBat \sep PCA pre-orthogonalisation \sep neuroscreening
\sep multi-centre study \sep machine learning \sep SHAP \sep DANN \sep
domain adaptation \sep logistic regression \sep two-classifier framework
\end{keyword}

\end{frontmatter}

% -----------------------------------------------------------------------
\section{Introduction}
\label{sec:introduction}

Electroencephalography (EEG) is a non-invasive, cost-effective
neurophysiological technique that holds substantial promise for
scalable, data-driven clinical diagnostics~\cite{poziomska2026quantity}.
The increasing availability of large, routinely collected clinical EEG
archives has opened new avenues for training machine-learning (ML)
classifiers capable of detecting neurological abnormalities at a
population scale~\cite{mushtaq2025eeg}.
A prerequisite for such systems to be deployable in practice is the
ability to generalise across acquisition environments: a model trained
on recordings from one set of hospitals must perform reliably on data
from institutions not represented during training.

This requirement is fundamentally challenged by the \emph{site
effect}---systematic, non-biological variance arising from
inter-centre differences in EEG hardware, electrode placement
conventions, amplifier settings, recording protocols, patient
preparation procedures, and patient
demography~\cite{jaramillo2024combat,li2022harmnqeeg}.
When ML models are trained on multi-centre data without addressing
the site effect, they can exploit institution-specific signatures in
the feature space rather than learning clinically discriminative
patterns, leading to marked degradation of performance when the model
encounters data from a previously unseen
hospital~\cite{kiessner2023extended,mellot2023meeg}.
As large, heterogeneous clinical EEG corpora become the natural
substrate for the next generation of diagnostic algorithms, the site
effect has emerged as one of the principal obstacles to building
generalisable neuroscreening models~\cite{dong2025normative,
mushtaq2025eeg}.

Statistical harmonisation methods, most prominently the ComBat
family~\cite{johnson2007adjusting,fortin2017harmonization,
fortin2018cortical}, were developed to address the analogous
batch-effect problem in genomics and neuroimaging, and have more
recently been applied to EEG
features~\cite{jaramillo2024combat,li2022harmnqeeg}.
CovBat~\cite{hu2023image} extends this framework to the multivariate
setting by additionally harmonising batch-specific structure in the
principal-component space of the residual covariance matrix.
Although computationally lightweight and widely available, all such
methods share a fundamental constraint: they operate feature by feature
and may therefore leave intact site-related structure encoded in
inter-feature covariances.

An alternative paradigm embeds harmonisation directly into the
representation-learning objective via Domain Adversarial Neural
Networks (DANNs)~\cite{ganin2016domain}, which have demonstrated
improved cross-subject and cross-site generalisation on EEG
classification tasks~\cite{bethge2022multidomain,liu2023multisource,
ozdenizci2020invariant}.
Critically, however, whether DANN-based harmonisation genuinely
suppresses the \emph{nonlinear} covariance component of the site
effect---or merely its linearly accessible component---has not been
systematically investigated.

The present study pursues two overarching goals.
First, we provide a rigorous characterisation of the site effect in a
corpus of 54\,779 EEG recordings from 30 Polish hospitals using a
\emph{two-classifier evaluation framework} that separates linearly
accessible site information from nonlinear covariance-based site
structure.
Second, we evaluate a comprehensive range of harmonisation strategies
under a unified experimental design, spanning simple site-wise
standardisation, neuroCombat, CovBat, PCA pre-orthogonalisation, and
four signal-level adversarial variants, and critically assess the
conditions under which these methods succeed or fail under both
standard cross-validation and a rigorous Leave-One-Site-Out (LOSO)
protocol.

The main contributions of this paper are as follows.
\begin{enumerate}
  \item We demonstrate that the site effect in a large clinical EEG
        corpus is strong and broadly distributed, with a nonlinear
        classifier achieving $\mathrm{MCC} = 0.865 \pm 0.003$ in a
        site identification task.
  \item We introduce a \emph{two-classifier evaluation framework}
        that uses a logistic regression linear probe and a CatBoost
        nonlinear probe applied symmetrically to every harmonisation
        condition, separating linearly accessible site information
        from nonlinear covariance structure.
  \item We reveal and characterise a \emph{harmonisation paradox}:
        marginal correction methods reduce the linear site-effect
        component (Cohen's~$d$, linear probe MCC near chance) yet
        simultaneously increase the nonlinear site-classification
        MCC towards a near-ceiling value, because equalising marginal
        distributions amplifies the salience of the surviving
        covariance-based site signatures.
  \item We identify, via SHAP analysis and ablation experiments, that
        inter-electrode coherence features are the primary carriers of
        site-specific information, with O1--O2 coherence and beta-band
        temporal coherence being particularly prominent.
  \item We show that PCA pre-orthogonalisation substantially
        attenuates the harmonisation paradox and stabilises
        cross-site neuroscreening performance across institutions.
  \item We evaluate four DANN variants under the same two-classifier
        framework, comparing signal-level adversarial harmonisation
        against all feature-level conditions on both site
        discriminability and downstream neuroscreening performance.
        Site discriminability of the adversarially learned
        representation is assessed via full-data MINet-DANN training
        with 5-fold cross-validation, whilst diagnostic performance
        is assessed exclusively under the LOSO protocol.
\end{enumerate}

% -----------------------------------------------------------------------
\section{Related Work}
\label{sec:related}

\subsection{The Site Effect in Multi-Centre EEG Studies}
\label{subsec:site_effect}

Systematic, non-biological variability arising from differences in
data acquisition conditions is a well-established problem across
quantitative biomedical sciences, where it is collectively termed the
\emph{batch effect}~\cite{johnson2007adjusting}.
In EEG research, the equivalent phenomenon is termed the \emph{site
effect}~\cite{li2022harmnqeeg}.
As large multi-centre EEG corpora become increasingly central to
data-driven clinical neuroscience~\cite{dong2025normative,
mushtaq2025eeg}, the site effect has emerged as a primary obstacle to
building classifiers that generalise beyond the hospitals represented
in training data.

Site effects arise from a heterogeneous combination of technical and
procedural factors including amplifier design, electrode materials,
impedance characteristics, recording-pipeline transfer functions,
activation procedures, resting-state recording duration, conventions
regarding eye-open versus eye-closed epochs, scalp preparation, and
patient alertness
management~\cite{jaramillo2024combat,mushtaq2025eeg}.
Demographic and clinical case-mix differences across institutions
further contribute to apparent site differences even when acquisition
protocols are nominally identical~\cite{dong2025normative}.
Because these sources act simultaneously and interact in complex ways,
the site effect is not a simple additive shift in individual features
but a high-dimensional, multivariate structure distributed across the
full feature space~\cite{mellot2023meeg}.

The degree to which site identity is encoded in EEG feature spaces
has been documented quantitatively.
Li~et~al.~\cite{li2022harmnqeeg} demonstrated systematic batch
effects across 12 devices and 9 countries in the HarMNqEEG corpus.
Jaramillo-Jimenez~et~al.~\cite{jaramillo2024combat} confirmed, using
$t$-SNE visualisation and Kruskal--Wallis tests, that spectral
features from five independent resting-state EEG datasets cluster by
site.
Kiessner~et~al.~\cite{kiessner2023extended} observed marked
performance degradation when a deep neural network trained on one
clinical corpus was evaluated on data from a different centre.
Mellot~et~al.~\cite{mellot2023meeg} showed that even
covariance-based EEG representations suffer substantial dataset shift
across acquisition devices and populations.

\subsection{Statistical Harmonisation Methods}
\label{subsec:statistical_harmonisation}

The canonical method is ComBat, originally proposed by
Johnson~et~al.~\cite{johnson2007adjusting} for batch-effect
correction in microarray gene-expression data.
Fortin~et~al.~\cite{fortin2017harmonization,fortin2018cortical}
established ComBat as the standard for retrospective neuroimaging
harmonisation.
CovBat~\cite{hu2023image} extends neuroCombat by additionally
harmonising the principal-component space of the residual covariance
matrix, thereby addressing batch effects residing in inter-feature
covariance structure.
CovBat has been shown to outperform standard ComBat in both
multivariate statistical evaluations and downstream machine-learning
metrics~\cite{hu2023image}.

In the EEG domain, Jaramillo-Jimenez~et~al.~\cite{jaramillo2024combat}
evaluated four ComBat variants on pooled resting-state EEG data from
five independent datasets and found that all variants reduced
batch-related dispersion in spectral features while preserving
age-related associations.
Li~et~al.~\cite{li2022harmnqeeg} demonstrated that ComBat-based
normalisation substantially reduced device-related clustering across
12 devices and 9 countries.

The most critical limitation of all ComBat-family methods is shared:
adjusting mean and variance independently for each feature leaves
intact any site-related structure residing in the \emph{covariance}
between features~\cite{hu2023image}.
The present study provides the first systematic empirical evidence for
this limitation in a large clinical EEG corpus through the
two-classifier evaluation framework.

\subsection{Deep Learning and Adversarial Harmonisation}
\label{subsec:dann_related}

An alternative paradigm embeds harmonisation directly into the
representation-learning objective, training a neural network to
produce features that are simultaneously predictive of the clinical
label and invariant to the recording site.
The canonical architecture is the Domain Adversarial Neural
Network~\cite{ganin2016domain}, which consists of a shared feature
extractor, a label predictor, and a domain discriminator connected via
a Gradient Reversal Layer (GRL).

Adversarial representation learning has been applied widely in the
EEG literature.
\"{O}zdenizci~et~al.~\cite{ozdenizci2020invariant} demonstrated
significant improvements in cross-subject generalisation on motor
imagery tasks.
Liu~et~al.~\cite{liu2023multisource} proposed MTDANN, which uses
EEGNet~\cite{lawhern2018eegnet} as the shared feature extractor and
outperformed single-source DANN and transfer learning baselines on
three public motor imagery datasets.
Bethge~et~al.~\cite{bethge2022multidomain} showed that a unified
adversarial invariance objective across multiple emotion recognition
datasets suppressed data-source-related information leakage by 35\%
while maintaining classification performance.

An important open question, which the present study directly
addresses, is whether adversarial training suppresses the
\emph{nonlinear covariance} component of the site effect or only its
linearly accessible component.
The two-classifier framework---applying both a linear probe and a
CatBoost nonlinear probe to the DANN representations---is designed to
provide a direct empirical answer to this question
(Section~\ref{subsubsec:repr_assessment}).

% -----------------------------------------------------------------------
\section{Materials and Methods}
\label{sec:methods}

\subsection{Dataset}
\label{subsec:dataset}

\subsubsection{The ELM19 Corpus and the
  \texorpdfstring{$\mathrm{ELM}_s$}{ELMs} Subset}

We use the ELM19 database~\cite{poziomska2026quantity}, a large-scale
multi-centre EEG corpus provided by Elmiko Biosignals sp.~z~o.o.\
(Milanówek, Poland), comprising 55\,787 clinical EEG recordings from
39 medical institutions across Poland during routine clinical
diagnostics.
Each recording is annotated with a binary normal/pathological label
derived from the accompanying EEG report following the procedure
detailed in~\cite{poziomska2026quantity}.
Patient ages range from 16 to 70 years; 52.9\% of recordings are
from female patients; and 52.2\% of recordings are labelled as
pathological.
EEG signals were acquired using the international 10--20 electrode
placement system with 19 channels, with sampling rates varying
between 200 and 500\,Hz across institutions.

Sites contributing fewer than 50 normal recordings were excluded to
ensure sufficient within-site representation for harmonisation
parameter estimation.
The resulting dataset---hereafter $\mathrm{ELM}_s$---consists of
54\,779 EEG recordings from 30 medical institutions, comprising
26\,568 normal (48.5\%) and 28\,211 pathological (51.5\%) recordings
from 28\,960 female subjects (52.9\%).

\subsubsection{Analysis Subset: Normal Recordings
  (\texorpdfstring{$\mathrm{ELM}_n$}{ELMn})}

To isolate the site effect from variance attributable to clinical
pathology, all primary analyses---exploratory data analysis, site
classification, harmonisation evaluation, and feature
importance---were conducted on the subset of recordings labelled as
normal, hereafter referred to as $\mathrm{ELM}_n$.
The full $\mathrm{ELM}_s$ dataset was used only in the downstream
neuroscreening experiment.

\subsection{EEG Preprocessing}
\label{subsec:preprocessing}

Raw EEG signals were preprocessed using
MNE-Python~\cite{gramfort2013mne} following the pipeline described
in~\cite{poziomska2026quantity}.
A Butterworth high-pass filter (cutoff 0.1\,Hz) removed slow drift,
and a low-pass filter (cutoff 40\,Hz) suppressed high-frequency
muscle artefacts.
A notch filter ($f_0 = 50$\,Hz, quality factor $Q = 5$) attenuated
power-line interference.
All signals were re-referenced to a common average reference and
resampled to 100\,Hz to ensure a uniform sampling rate across all
institutions.
Continuous recordings were segmented into non-overlapping 6-second
frames (600 samples at 100\,Hz).
Frames containing flat-line channels or voltages exceeding
$800\,\mu\text{V}$ were discarded as artefacts.

\subsection{Feature Extraction}
\label{subsec:features}

Three groups of features were extracted from each 6-second frame and
subsequently aggregated per recording using the median across all
valid frames~\cite{poziomska2026quantity}.

\subsubsection{Riemannian Covariance Features}

For each frame, the $19 \times 19$ channel covariance matrix was
computed and averaged using the Riemannian geometric mean via the
\texttt{pyRiemann} library.
The averaged matrix was projected onto its tangent space to obtain a
Euclidean representation; the lower triangular elements were
extracted, yielding \textbf{190 features} per recording.

\subsubsection{Multitaper Power Features}

Power spectral densities were estimated using the multitaper
method~\cite{thomson1982spectrum} and integrated across 14 custom
overlapping frequency bands covering 0.5--40\,Hz, spanning the
standard delta, theta, alpha, and beta EEG bands.
PSDs were normalised such that the sum across all channels and bands
equalled one within each frame, yielding \textbf{266 power features}
per recording.

\subsubsection{Inter-Electrode Coherence Features}

Band-wise magnitude-squared coherence was computed for all
$\binom{19}{2} = 171$ inter-electrode pairs in each of the 14
frequency bands:
\begin{equation}
  C_{xy}(f_b) = \frac{\bigl|S_{xy}(f_b)\bigr|^2}
                     {S_{xx}(f_b)\cdot S_{yy}(f_b)},
  \label{eq:coherence}
\end{equation}
where $S_{xy}(f_b)$ is the cross-spectral density and $S_{xx}(f_b)$,
$S_{yy}(f_b)$ are the auto-spectral densities of channels $x$ and $y$
in band $f_b$, yielding $171 \times 14 = \mathbf{2\,394}$
\textbf{coherence features} per recording.

\subsubsection{Final Feature Vector}

Concatenating the three groups yields a final feature vector of
$190 + 266 + 2\,394 = \mathbf{2\,850}$ features per recording.

\subsection{Approach~1: Feature-Level Harmonisation}
\label{subsec:approach1}

All feature-level methods were applied independently to each of the
2\,850 features.
Harmonisation parameters were estimated exclusively on the training
portion of each data split and subsequently applied to both training
and test data to prevent data leakage.

\subsubsection{Site-Wise Z-Score Standardisation}
\label{subsubsec:sitewise}

As a conceptually simple baseline, each feature was standardised to
zero mean and unit variance using within-site statistics estimated
from the training split:
\begin{equation}
  z_{ij} = \frac{x_{ij} - \hat{\mu}_{s(i),j}}
                {\hat{\sigma}_{s(i),j}},
  \label{eq:sitewise}
\end{equation}
where $\hat{\mu}_{s(i),j}$ and $\hat{\sigma}_{s(i),j}$ are the
within-site mean and standard deviation of feature $j$ estimated on
the training data.
This transformation equalises the first two moments of each feature
distribution across sites but makes no assumptions about the structure
of the site effect, does not account for biological covariates, and
leaves inter-feature covariance structure entirely uncorrected.

\subsubsection{neuroCombat}
\label{subsubsec:neurocombat}

The neuroCombat implementation of the ComBat
algorithm~\cite{johnson2007adjusting,fortin2017harmonization} was
applied with the hospital identifier as the batch variable and patient
age and sex as biological covariates.
ComBat models the observed feature value $y_{ij}$ for subject $i$ and
feature $j$ as
\begin{equation}
  y_{ij} = \alpha_j + \mathbf{x}_i^{\top}\boldsymbol{\beta}_j
           + \gamma_{s(i),j} + \delta_{s(i),j}\,\varepsilon_{ij},
  \label{eq:combat}
\end{equation}
where $\alpha_j$ is the global intercept, $\mathbf{x}_i$ is the
vector of biological covariates with coefficients
$\boldsymbol{\beta}_j$, $\gamma_{s(i),j}$ is an additive site
effect, $\delta_{s(i),j}$ is a multiplicative site effect, and
$\varepsilon_{ij}$ is residual noise, all estimated via empirical
Bayes~\cite{johnson2007adjusting}.
Harmonised values are recovered by:
\begin{equation}
  y^*_{ij} = \frac{y_{ij} - \hat{\alpha}_j
             - \mathbf{x}_i^{\top}\hat{\boldsymbol{\beta}}_j
             - \hat{\gamma}_{s(i),j}}{\hat{\delta}_{s(i),j}}
             + \hat{\alpha}_j
             + \mathbf{x}_i^{\top}\hat{\boldsymbol{\beta}}_j.
  \label{eq:combat_corrected}
\end{equation}

\subsubsection{CovBat}
\label{subsubsec:covbat}

CovBat~\cite{hu2023image} first applies standard ComBat
(Equations~\ref{eq:combat} and~\ref{eq:combat_corrected}) to correct
marginal mean and variance shifts, then performs a second
harmonisation step in the principal-component space of the residual
covariance matrix, targeting site-specific shifts in this PC space.
We used the implementation of Hu~et~al.~\cite{hu2023image} with
patient age and sex as biological covariates.

\subsubsection{PCA Pre-Orthogonalisation}
\label{subsubsec:pca}

As an additional preprocessing step---applied \emph{before} any of
the marginal correction methods---we evaluated PCA
pre-orthogonalisation of the feature space.
A PCA projection was computed on the training data (retaining 100\%
of variance) and applied to both training and test data prior to
harmonisation.
Because the resulting PC features are mutually uncorrelated by
construction, marginal correction methods subsequently applied to
this transformed space cannot exploit or amplify inter-feature
covariance structure.
PCA pre-orthogonalisation was evaluated both as a standalone
condition and in combination with site-wise standardisation,
neuroCombat, and CovBat.

\subsubsection{Downstream Classifier for Neuroscreening}
\label{subsubsec:classifier}

Original or harmonised feature vectors were passed to a
gradient-boosted ensemble (GBE) model composed of 30 individual
CatBoost~\cite{prokhorenkova2018catboost} classifiers, identical to
the architecture described in~\cite{poziomska2026quantity}.
Hyperparameters were fixed to the values reported
in~\cite{poziomska2026quantity} to ensure comparability.

\subsection{Approach~2: Signal-Level Adversarial Harmonisation}
\label{subsec:approach2_dann}

Feature-level harmonisation methods operate on pre-extracted,
hand-crafted representations and, as will be demonstrated in
Section~\ref{subsubsec:paradox}, cannot address site-related structure
residing in inter-feature covariances or their nonlinear dependencies.
To investigate whether this limitation can be overcome through
end-to-end representation learning, we develop a signal-level approach
based on Domain Adversarial Neural Networks~\cite{ganin2016domain},
integrated with the multiple-instance learning (MIL) architecture
of~\cite{poziomska2026quantity}.

Crucially, the adversarial objective requires the network to solve
both the pathology classification task and the site discrimination task
simultaneously.
Because the pathology labels are required for the classification head,
adversarial training is necessarily conducted on the full
$\mathrm{ELM}_s$ dataset rather than on the normal-only
$\mathrm{ELM}_n$ subset used in the feature-level experiments.

\subsubsection{Base Architecture}
\label{subsubsec:dann_arch}

The base architecture is the MINet model
of~\cite{poziomska2026quantity}, consisting of three components
(Figure~\ref{fig:dann_arch}).

\paragraph{Frame encoder.}
An EEGNet-based~\cite{lawhern2018eegnet} convolutional feature
extractor $G_f(\cdot\,;\,\theta_f)$ maps each 6-second EEG frame
$\mathbf{x} \in \mathbb{R}^{C \times T}$ (19 channels, 600 time
points at 100\,Hz) to a 288-dimensional frame-level latent
representation $\mathbf{f} \in \mathbb{R}^{288}$.
The encoder comprises 1\,408 parameters and performs separate
temporal and spatial convolutions following~\cite{lawhern2018eegnet}.

\paragraph{Attention-based recording aggregation.}
Frame-level representations are aggregated using the non-gated
attention module of~\cite{poziomska2026quantity}, which incorporates
query--key--value structure borrowed from the transformer
architecture.
The module assigns a scalar weight $a_t$ to each frame and computes
the recording-level representation as
\begin{equation}
  \mathbf{h} = \sum_{t=1}^{T_r} a_t\,\mathbf{f}_t,
  \qquad
  a_t = \frac{\exp\!\bigl(\mathbf{w}^\top
              \tanh(\mathbf{V}\mathbf{f}_t)\bigr)}
             {\sum_{\tau=1}^{T_r}
              \exp\!\bigl(\mathbf{w}^\top
              \tanh(\mathbf{V}\mathbf{f}_\tau)\bigr)},
  \label{eq:attention}
\end{equation}
where $T_r$ is the number of valid frames in the recording and
$\mathbf{V}$, $\mathbf{w}$ are learned parameters.

\paragraph{Pathology classifier.}
The recording-level representation $\mathbf{h}$ is passed to a
pathology classifier $G_y(\cdot\,;\,\theta_y)$---a two-layer fully
connected network producing a binary normal/pathological prediction.

\paragraph{Site discriminator.}
A separate fully connected network $G_d(\cdot\,;\,\theta_d)$ maps
$\mathbf{h}$ to a distribution over the 30 recording institutions.
Its connection to the backbone differs across the four training
variants described in Section~\ref{subsubsec:dann_variants}: in
MINet-DANN it is connected via a Gradient Reversal Layer
(GRL)~\cite{ganin2016domain}, which leaves the input unchanged during
forward propagation and multiplies the gradient by $-\lambda$ during
back-propagation; in MINet-MTL it is connected via an identity
mapping so that gradients flow unmodified; and in MINet no site
discriminator is attached.

\begin{figure}[htbp]
  \centering
  \includegraphics[width=0.85\textwidth]{Figs/dann_arch.png}
  \caption{Architecture of the Domain Adversarial Neural Network
    used for signal-level harmonisation.
    A shared feature extractor feeds both a pathology classifier and
    a site discriminator; the Gradient Reversal Layer (GRL) inverts
    the sign of gradients from the discriminator during
    back-propagation, training the extractor to produce
    site-invariant representations.}
  \label{fig:dann_arch}
\end{figure}

\subsubsection{Training Variants}
\label{subsubsec:dann_variants}

Four training conditions are evaluated to isolate the contribution
of each component of the adversarial framework.
All four conditions are trained under the same Leave-One-Site-Out
(LOSO) protocol as the feature-level experiments
(Section~\ref{subsubsec:clinical_eval}): in each round, all
recordings from one institution form the test set and the model is
trained on the remaining 29 institutions (source pool).

\paragraph{MINet (signal-level baseline).}
The MINet baseline uses the pretrained backbone
of~\cite{poziomska2026quantity} without any domain objective
($\lambda \equiv 0$).
The encoder is first pretrained on individual frames in a
single-instance fashion identical to the siNet pretraining
of~\cite{poziomska2026quantity}, after which the full MINet network
is trained for pathology classification on the source pool.
No site discriminator is attached.
This condition establishes the degree of site information retained
in the learned representation $\mathbf{h}$ in the absence of any
adversarial pressure.

\paragraph{MINet-MTL (multi-task baseline, no gradient reversal).}
MINet-MTL replaces the GRL with an identity mapping, so that the
gradient from the site discriminator $G_d$ flows unmodified into the
backbone.
The network is initialised with random weights and trained from
scratch for 50 epochs.
The combined loss is:
\begin{equation}
  \mathcal{L}_{\mathrm{MTL}} =
    5.0 \cdot \mathcal{L}_y\!\bigl(G_y(\mathbf{h}),\,y\bigr)
    + \mathcal{L}_d\!\bigl(G_d(\mathbf{h}),\,d\bigr),
  \label{eq:mtl_loss}
\end{equation}
where $y$ is the pathology label and $d$ is the site label.
Because the site loss flows normally into the backbone, MINet-MTL
is expected to \emph{exploit} site identity rather than suppress it,
producing a larger CatBoost-MCC on $\mathbf{h}$ than MINet and
thereby providing an empirical upper bound on the site information
that an unconstrained multi-task network would encode in the
recording-level representation.

\paragraph{MINet-DANN (full adversarial training).}
MINet-DANN inserts the GRL between the attention module and the site
discriminator, replacing the site loss term with an adversarial
objective.
The network is initialised with random weights and trained from
scratch for 50 epochs.
The saddle-point objective is:
\begin{equation}
  \mathcal{L}(\theta_f,\theta_a,\theta_y,\theta_d) =
      5.0 \cdot \mathcal{L}_y\!\bigl(G_y(\mathbf{h}),\,y\bigr)
      - \lambda\,\mathcal{L}_d\!\bigl(G_d(\mathbf{h}),\,d\bigr),
  \label{eq:dann_loss}
\end{equation}
where $\theta_f$, $\theta_a$, $\theta_y$, and $\theta_d$ are the
parameters of the frame encoder, attention module, pathology
classifier, and site discriminator, respectively.
The domain-adaptation weight $\lambda$ is annealed from 0 to 1
during training according to the schedule of
Ganin~et~al.~\cite{ganin2016domain}:
\begin{equation}
  \lambda_p = \frac{2}{1 + \exp(-\gamma p)} - 1,
  \qquad \gamma = 10,
  \label{eq:lambda_schedule}
\end{equation}
where $p \in [0,1]$ denotes training progress (fraction of total
epochs completed).
Model selection at each LOSO round is based on the highest
pathology AUC obtained on a held-out validation subset drawn from
the source pool.

\paragraph{MINet-DANN+tuning (frozen backbone, retrained classifier).}
MINet-DANN+tuning takes the best MINet-DANN checkpoint from each
LOSO round, freezes the parameters of the frame encoder $G_f$,
attention module, and site discriminator $G_d$, and retrains only
the pathology classifier $G_y$ on the source-pool data using the
AdamW optimiser for 30 epochs ($\lambda \equiv 0$ throughout).
This recalibration step adjusts the diagnostic decision boundaries
within the site-agnostic latent space learned by MINet-DANN, without
modifying the harmonised representation.
Because the target site is never exposed to the model during this
step, the condition remains strictly LOSO-compliant and provides a
calibrated diagnostic estimate directly comparable to the
feature-level LOSO conditions.

% -----------------------------------------------------------------------
\subsection{Evaluation Framework}
\label{subsec:evaluation}

\subsubsection{Two-Classifier Site-Discriminability Evaluation}
\label{subsubsec:two_classifier}

The site effect was quantified using a supervised site-classification
task on $\mathrm{ELM}_n$, using 5-fold stratified cross-validation
with stratification by hospital~ID.
To distinguish between the linearly accessible and nonlinearly
accessible components of the site effect, two classifiers of
contrasting capacity were applied to every condition:

\begin{itemize}
  \item \textbf{Logistic regression (LR)} with $\ell_2$
        regularisation, selected by inner cross-validation on the
        training folds. LR can only exploit mean shifts and linear
        inter-feature relationships, and therefore measures the
        component of the site effect that survives as a linearly
        separable signal.
  \item \textbf{CatBoost gradient-boosted classifier}, which can
        additionally exploit higher-order inter-feature interactions
        and covariance structure, and therefore measures the
        \emph{total} site information including structure that is
        invisible to a linear classifier.
\end{itemize}

The gap between LR-MCC and CatBoost-MCC under any given condition
directly quantifies the degree to which that condition's site effect
is encoded in nonlinear covariance structure.
For feature-level conditions, both classifiers are trained on the
2\,850-dimensional feature vector.
For signal-level conditions, both classifiers are trained on the
recording-level representation $\mathbf{h}$ extracted from
full-data MINet-DANN training, as described in
Section~\ref{subsubsec:repr_assessment}.
The Matthews Correlation Coefficient (MCC) is used as the primary
metric:
\begin{equation}
  \mathrm{MCC} = \frac{c\,s - \sum_{k=1}^{K} p_k t_k}
    {\sqrt{\!\left(s^2 - \sum_{k=1}^{K} p_k^2\right)
           \!\left(s^2 - \sum_{k=1}^{K} t_k^2\right)}},
  \label{eq:mcc}
\end{equation}
where $c = \sum_k C_{kk}$, $s = \sum_{i,j} C_{ij}$,
$t_k = \sum_j C_{kj}$, and $p_k = \sum_i C_{ik}$ for confusion
matrix $\mathbf{C} \in \mathbb{N}^{K \times K}$, and $K$ denotes
the number of classes ($K = 30$ for the site-classification task;
$K = 2$ for the binary neuroscreening task).

\subsubsection{Representation Quality Assessment for Signal-Level
  Models}
\label{subsubsec:repr_assessment}

To provide a direct comparison between the site information encoded
in the DANN representations and that suppressed by the statistical
harmonisation methods, an additional assessment is performed on the
representations learned by the full-data DANN models.
Specifically, MINet-DANN is trained from scratch on the complete
$\mathrm{ELM}_s$ dataset (all 30 sites, without held-out folds),
and the resulting recording-level representations $\mathbf{h}$ are
extracted for all recordings.
Both LR and CatBoost classifiers are then trained to predict the
site label from $\mathbf{h}$ using 5-fold stratified
cross-validation, stratified by hospital~ID, and performance is
reported as MCC (mean $\pm$ SD across folds).

This condition directly parallels the feature-level site-classification
experiments---in which LR and CatBoost were trained on the original
and harmonised 2\,850-dimensional feature vectors---and yields a
meaningful comparison of the residual site information in:
\begin{enumerate}
  \item the raw feature space,
  \item each statistical harmonisation condition, and
  \item the adversarially learned representation $\mathbf{h}$.
\end{enumerate}
Because DANN training on the full dataset maximises the signal
available to the adversarial objective, the MCC values from this
condition represent the strongest achievable site-invariance of the
DANN representation and serve as an upper bound on the harmonisation
effectiveness of the signal-level approach.

\subsubsection{Feature Importance Analysis}
\label{subsubsec:shap}

SHAP (SHapley Additive exPlanations) values were computed for the
trained CatBoost site-classification model using the SHAP
TreeExplainer:
\begin{equation}
  g(\mathbf{z}') = \phi_0 + \sum_{i=1}^{M} \phi_i z'_i,
  \label{eq:shap}
\end{equation}
where $\phi_0$ is the base value and $\phi_i$ is the SHAP value for
feature $i$.
Global feature importance was summarised by the mean absolute SHAP
value across all predictions and classes.
An ablation study trained models on each feature group individually
(covariance, power, coherence) and with each group sequentially
removed.

\subsubsection{Univariate Harmonisation Effectiveness: Cohen's~$d$}
\label{subsubsec:cohend}

The impact of harmonisation on individual feature distributions was
quantified using the standardised mean difference between
institutions:
\begin{equation}
  d = \frac{\mu_1 - \mu_2}{\sigma_{\mathrm{pooled}}},
  \qquad
  \sigma_{\mathrm{pooled}} =
    \sqrt{\frac{(n_1-1)s_1^2 + (n_2-1)s_2^2}{n_1+n_2-2}},
  \label{eq:cohend}
\end{equation}
computed for each feature over every pair of institutions and averaged
across features to produce an all-vs-all pairwise summary.

\subsubsection{Clinical Neuroscreening Evaluation}
\label{subsubsec:clinical_eval}

Clinical utility was assessed on the binary pathology classification
task using the full $\mathrm{ELM}_s$ dataset under two protocols.

\paragraph{5-Fold Cross-Validation.}
The dataset was split into five folds stratified by hospital~ID,
pathology label, and sex, and performance (MCC and AUC) was averaged
across folds.
This protocol provides a general performance estimate representative
of the setting in which all sites are seen during training.
This protocol is applied exclusively to the feature-level conditions
(Approach~1); the DANN variants (Approach~2) are evaluated under the
LOSO protocol only (see below).

\paragraph{Leave-One-Site-Out (LOSO).}
In each round, all recordings from one of the 30 institutions are
held out as the target test set, while the remaining 29 institutions
form the training source pool.
A calibration sample of $n_{\mathrm{cal}} = 30$ normal recordings
from the held-out site is reserved prior to each round.
For feature-level methods, this calibration sample is used together
with the training data from all other institutions to fit the
harmonisation parameters; for the unharmonised feature-level baseline
the calibration sample is discarded.
For signal-level DANN variants, the calibration sample is included in
the source-pool training data.
Because $n_{\mathrm{cal}} = 30$ represents a negligible fraction of
the total training corpus (fewer than 0.06\% of all recordings), this
is not expected to introduce meaningful bias into the evaluation.
The calibration sample is never included in the pathology
classification training set under any condition.
This protocol directly simulates deployment to a previously unseen
recording environment and is the primary measure of generalisation in
the present study.

Notably, for the deep learning-based domain adaptation variants
(Approach~2: MINet, MINet-MTL, MINet-DANN, and MINet-DANN+tuning),
the LOSO protocol is the \emph{exclusive} method of evaluation.
The adversarial framework is explicitly designed to maximise
generalisation to novel clinical institutions without prior exposure to
their signals; the strictly enforced LOSO evaluation therefore
provides the most relevant assessment of these models' deployment
potential.
Performance (MCC and AUC) is reported both as an aggregate across
all 30 LOSO rounds and as a per-institution profile.

% -----------------------------------------------------------------------
\section{Results}
\label{sec:results}

\subsection{Exploratory Analysis of Inter-Institutional Heterogeneity}
\label{subsec:exploratory}

Exploratory analysis of the $\mathrm{ELM}_n$ subset revealed
substantial heterogeneity across the 30 institutions in patient
demographics and recording protocols.
Age distributions varied markedly between sites: institutions such as
PIO, SL2, and TER\_L served predominantly young populations (median
ages 17.6, 19.5, and 17.8 years, respectively), while institutions
such as MOR and TOR served predominantly older patients (median ages
52.1 and 50.8 years, respectively).
The maximum difference in median age between any two institutions
exceeded 30 years (full statistics are given in
Table~\ref{tab:age_stats} in the Appendix).
Recording duration showed further heterogeneity: institution
KATMOJPRZ recorded sessions lasting up to 8\,322 seconds, and ZOZLO
showed a bimodal distribution suggesting the use of multiple recording
protocols within a single institution.

\subsection{Quantifying the Site Effect}
\label{subsec:site_classification}

\subsubsection{Site Classification Performance}

The CatBoost nonlinear classifier trained on the unharmonised
$\mathrm{ELM}_n$ feature set achieved an overall MCC of
$0.865 \pm 0.003$ across the 30 institutions in 5-fold stratified
cross-validation, confirming that site identity is strongly and
reliably encoded in the feature space.
The logistic regression linear probe achieved an overall MCC of
$\mathbf{[TBF]} \pm \mathbf{[TBF]}$ on the same unharmonised data.
The LR--CatBoost gap of $\mathbf{[TBF]}$ on the unharmonised baseline
establishes the degree to which the site effect in the raw feature
space is already encoded in nonlinear covariance structure.

Per-institution one-vs-rest MCC scores for the CatBoost classifier
are summarised in Table~\ref{tab:per_site_mcc}.
A total of 23 out of 30 institutions achieved average MCC scores
above 0.7.
A small subset of institutions---AKS, CHE, KIEG, and
SL2---performed poorly (MCC 0.39--0.52), likely owing to
distributional similarity with neighbouring sites in the corpus.
Institution CMD achieved the lowest average MCC of 0.21.

\begin{table}[htbp]
  \centering
  \caption{Per-institution one-vs-rest MCC (mean $\pm$ SD, 5-fold CV
    on $\mathrm{ELM}_n$, unharmonised CatBoost classifier).
    Institutions are ordered by decreasing MCC, reading the left
    column top-to-bottom then the right column top-to-bottom.}
  \label{tab:per_site_mcc}
  \begin{tabular}{lc@{\quad}lc}
    \toprule
    \textbf{Institution} & \textbf{MCC} &
    \textbf{Institution} & \textbf{MCC} \\
    \midrule
    MOR       & $0.954 \pm 0.035$ & PRZ      & $0.844 \pm 0.016$ \\
    B2K       & $0.952 \pm 0.006$ & KLU      & $0.828 \pm 0.015$ \\
    TOR       & $0.938 \pm 0.015$ & WLU      & $0.805 \pm 0.015$ \\
    KAL       & $0.941 \pm 0.030$ & MKW      & $0.809 \pm 0.034$ \\
    LUMICE    & $0.947 \pm 0.034$ & SRK      & $0.780 \pm 0.036$ \\
    KUD       & $0.933 \pm 0.017$ & OST      & $0.781 \pm 0.016$ \\
    KATMOJPRZ & $0.918 \pm 0.012$ & ZOZLO    & $0.782 \pm 0.048$ \\
    STG1      & $0.922 \pm 0.010$ & OTW      & $0.830 \pm 0.010$ \\
    Z04O      & $0.915 \pm 0.012$ & GAK      & $0.748 \pm 0.014$ \\
    SZC       & $0.893 \pm 0.005$ & TER\_L   & $0.682 \pm 0.025$ \\
    SLU       & $0.890 \pm 0.026$ & PIO      & $0.626 \pm 0.098$ \\
    PUS       & $0.899 \pm 0.048$ & AKS      & $0.517 \pm 0.117$ \\
    ARCHDAM   & $0.872 \pm 0.014$ & SL2      & $0.467 \pm 0.112$ \\
    LUX\_A    & $0.851 \pm 0.016$ & CHE      & $0.407 \pm 0.113$ \\
    overall   & $0.865 \pm 0.003$ & KIEG     & $0.390 \pm 0.079$ \\
              &                   & CMD      & $0.207 \pm 0.047$ \\
    \bottomrule
  \end{tabular}
\end{table}

\subsubsection{Feature Importance: SHAP Analysis and Ablation}
\label{subsubsec:shap_results}

Table~\ref{tab:ablation} summarises site-classification MCC for models
trained on individual feature groups and with each group excluded.
A model trained exclusively on coherence features achieved MCC of
$0.850 \pm 0.003$, nearly matching the full-feature-set performance.
Power features alone achieved $0.600 \pm 0.005$, and covariance
features alone achieved $0.752 \pm 0.009$.
Excluding coherence features reduced performance most substantially
(to $0.785 \pm 0.007$), while excluding power or covariance features
had comparatively little effect, confirming that coherence features
are the dominant carriers of site-specific information.

The SHAP analysis corroborates these findings.
Among the coherence features, the O1--O2 electrode pair achieved the
highest importance scores across the majority of frequency bands.
Beta-band coherence was particularly prominent in the temporal regions
(T3, T4, T5, T6, F7, F8).
Among spectral power features, delta-band power and the power at
electrode Fz were the most important.
Covariance features contributed comparatively little to the model's
site predictions.

\begin{table}[htbp]
  \centering
  \caption{Overall site-classification MCC (mean $\pm$ SD, 5-fold CV
    on $\mathrm{ELM}_n$, CatBoost classifier) for the feature-group
    ablation study.}
  \label{tab:ablation}
  \begin{tabular}{lcc}
    \toprule
    \textbf{Condition} & \textbf{MCC} & \textbf{SD} \\
    \midrule
    \multicolumn{3}{l}{\textit{Single feature group}} \\
    \quad Coherence only  & 0.850 & 0.003 \\
    \quad Power only      & 0.600 & 0.005 \\
    \quad Covariance only & 0.752 & 0.009 \\
    \midrule
    \multicolumn{3}{l}{\textit{One group excluded}} \\
    \quad All except coherence  & 0.785 & 0.007 \\
    \quad All except power      & 0.857 & 0.002 \\
    \quad All except covariance & 0.863 & 0.004 \\
    \bottomrule
  \end{tabular}
\end{table}

\subsection{The Two-Classifier Evaluation of Harmonisation Methods}
\label{subsec:harmonisation_results}

\subsubsection{Cohen's~$d$: Univariate Harmonisation Effectiveness}

Prior to harmonisation, substantial inter-site differences were
present at the feature level, with mean pairwise Cohen's~$d$
of $\mathbf{[TBF]}$ averaged across all institution pairs and all
2\,850 features.
All marginal correction methods reduced inter-site Cohen's~$d$
relative to the unharmonised baseline (Table~\ref{tab:cohend_summary}).

\begin{table}[htbp]
  \centering
  \caption{Mean pairwise Cohen's~$d$ (averaged across all institution
    pairs and all 2\,850 features) before and after each harmonisation
    condition, applied to $\mathrm{ELM}_n$.}
  \label{tab:cohend_summary}
  \begin{tabular}{lc}
    \toprule
    \textbf{Condition} & \textbf{Mean pairwise Cohen's~$d$} \\
    \midrule
    Original (unharmonised)         & \textbf{[TBF]} \\
    PCA pre-orthogonalisation only  & \textbf{[TBF]} \\
    Site-wise standardisation       & \textbf{[TBF]} \\
    PCA + site-wise stand.          & \textbf{[TBF]} \\
    neuroCombat                     & \textbf{[TBF]} \\
    PCA + neuroCombat               & \textbf{[TBF]} \\
    CovBat                          & \textbf{[TBF]} \\
    PCA + CovBat                    & \textbf{[TBF]} \\
    \bottomrule
  \end{tabular}
\end{table}

\subsubsection{The Harmonisation Paradox Revealed by the
  Two-Classifier Framework}
\label{subsubsec:paradox}

Table~\ref{tab:site_class_harmonisation} reports the LR-MCC and
CatBoost-MCC for all conditions.
The logistic regression linear probe assigned near-chance MCC to all
harmonised conditions, confirming that marginal correction
successfully removes the \emph{linear} component of the site effect.
The CatBoost nonlinear probe, by contrast, assigned near-ceiling MCC
($> 0.98$) to all harmonised conditions applied without PCA
pre-orthogonalisation, demonstrating the \emph{harmonisation paradox}:
the nonlinear covariance component of the site effect not only
survives marginal correction but becomes \emph{more} salient to the
nonlinear classifier.

The LR--CatBoost gap is therefore the quantitative fingerprint of
the harmonisation paradox: a large gap indicates that the residual
site effect after correction is encoded almost entirely in covariance
structure.
PCA pre-orthogonalisation substantially reduces this gap by
destroying the inter-feature covariance prior to correction,
attenuating the paradox across all harmonisation conditions.

The signal-level conditions (MINet, MINet-MTL, MINet-DANN,
MINet-DANN+tuning) are evaluated via the same two-classifier framework
applied to the representation $\mathbf{h}$.
\textbf{[TBF: report LR-MCC and CatBoost-MCC for all four DANN
variants once experiments are complete, and comment on whether
MINet-DANN achieves a smaller LR--CatBoost gap than the
feature-level conditions.]}

Figure~\ref{fig:site_mcc_no_pca} shows per-institution
site-classification MCC (CatBoost) without PCA
pre-orthogonalisation, and Figure~\ref{fig:site_mcc_pca} shows the
analogous results with PCA pre-orthogonalisation.

\begin{figure}[htbp]
  \centering
  \includegraphics[width=\textwidth]{Figs/site_mcc_pca_00.pdf}
  \caption{Per-institution site-classification MCC (CatBoost,
    one-vs-rest, 5-fold CV on $\mathrm{ELM}_n$) for the raw feature
    set and after four harmonisation methods \emph{without} PCA
    pre-orthogonalisation.
    All four methods drive MCC towards a near-ceiling value,
    demonstrating the harmonisation paradox.}
  \label{fig:site_mcc_no_pca}
\end{figure}

\begin{figure}[htbp]
  \centering
  \includegraphics[width=\textwidth]{Figs/site_mcc_pca_10.pdf}
  \caption{Per-institution site-classification MCC (CatBoost,
    one-vs-rest, 5-fold CV on $\mathrm{ELM}_n$) after four
    harmonisation methods \emph{with} full-variance PCA
    pre-orthogonalisation.
    PCA pre-orthogonalisation substantially attenuates the
    harmonisation paradox.}
  \label{fig:site_mcc_pca}
\end{figure}

\begin{table}[htbp]
  \centering
  \caption{Overall site-classification MCC (mean $\pm$ SD, 5-fold CV
    on $\mathrm{ELM}_n$) for all harmonisation conditions under the
    two-classifier framework.
    LR = logistic regression linear probe; CB = CatBoost nonlinear
    probe.
    The LR--CB gap quantifies the degree to which residual site
    information is encoded in nonlinear covariance structure.
    $^\dagger$Signal-level models are evaluated under the LOSO
    protocol only (Section~\ref{subsubsec:clinical_eval}); site
    discriminability is assessed on representations from full-data
    training (Section~\ref{subsubsec:repr_assessment}).}
  \label{tab:site_class_harmonisation}
  \begin{tabular}{lcccc}
    \toprule
    & \multicolumn{2}{c}{\textbf{LR (linear)}}
    & \multicolumn{2}{c}{\textbf{CatBoost (nonlinear)}} \\
    \cmidrule(lr){2-3}\cmidrule(lr){4-5}
    \textbf{Condition} & MCC & SD & MCC & SD \\
    \midrule
    \multicolumn{5}{l}{\textit{Feature-level, without PCA}} \\
    \quad Original (unharmonised)
      & \textbf{[TBF]} & \textbf{[TBF]} & 0.865 & 0.003 \\
    \quad Site-wise standardisation
      & \textbf{[TBF]} & \textbf{[TBF]} & 0.986 & 0.003 \\
    \quad neuroCombat
      & \textbf{[TBF]} & \textbf{[TBF]} & 0.983 & 0.002 \\
    \quad CovBat
      & \textbf{[TBF]} & \textbf{[TBF]} & \textbf{[TBF]} & \textbf{[TBF]} \\
    \midrule
    \multicolumn{5}{l}{\textit{Feature-level, with PCA pre-orthogonalisation}} \\
    \quad PCA only
      & \textbf{[TBF]} & \textbf{[TBF]} & \textbf{[TBF]} & \textbf{[TBF]} \\
    \quad PCA + site-wise stand.
      & \textbf{[TBF]} & \textbf{[TBF]} & \textbf{[TBF]} & \textbf{[TBF]} \\
    \quad PCA + neuroCombat
      & \textbf{[TBF]} & \textbf{[TBF]} & \textbf{[TBF]} & \textbf{[TBF]} \\
    \quad PCA + CovBat
      & \textbf{[TBF]} & \textbf{[TBF]} & \textbf{[TBF]} & \textbf{[TBF]} \\
    \midrule
    \multicolumn{5}{l}{\textit{Signal-level (both probes on $\mathbf{h}$)$^\dagger$}} \\
    \quad MINet ($\lambda \equiv 0$)
      & \textbf{[TBF]} & \textbf{[TBF]} & \textbf{[TBF]} & \textbf{[TBF]} \\
    \quad MINet-MTL
      & \textbf{[TBF]} & \textbf{[TBF]} & \textbf{[TBF]} & \textbf{[TBF]} \\
    \quad MINet-DANN
      & \textbf{[TBF]} & \textbf{[TBF]} & \textbf{[TBF]} & \textbf{[TBF]} \\
    \quad MINet-DANN+tuning
      & \textbf{[TBF]} & \textbf{[TBF]} & \textbf{[TBF]} & \textbf{[TBF]} \\
    \bottomrule
  \end{tabular}
\end{table}

\subsection{Impact of Harmonisation on Clinical Neuroscreening}
\label{subsec:clinical_results}

\subsubsection{Per-Institution LOSO Performance}

Figure~\ref{fig:patho_auc_no_pca} shows the per-institution
neuroscreening AUC under the LOSO protocol for all feature-level
harmonisation methods applied without PCA pre-orthogonalisation.
Without PCA, the four methods diverge noticeably at individual
institutions.
\textbf{[TBF: identify specific outlier institutions, describe the
direction of the anomaly, and comment on likely causes.]}

Figure~\ref{fig:patho_auc_pca} shows the analogous results with PCA
pre-orthogonalisation.
In contrast to the no-PCA condition, the four methods cluster closely
together across all institutions, and the institution-level
variability in AUC is markedly reduced.

\begin{figure}[htbp]
  \centering
  \includegraphics[width=\textwidth]{Figs/patho_auc_pca_none.pdf}
  \caption{Per-institution neuroscreening AUC under the LOSO
    protocol for the GBE classifier applied to raw features and
    after four feature-level harmonisation methods, \emph{without}
    PCA pre-orthogonalisation.}
  \label{fig:patho_auc_no_pca}
\end{figure}

\begin{figure}[htbp]
  \centering
  \includegraphics[width=\textwidth]{Figs/patho_auc_pca_all.pdf}
  \caption{Per-institution neuroscreening AUC under the LOSO
    protocol for the GBE classifier applied to raw features and
    after four feature-level harmonisation methods, \emph{with}
    full-variance PCA pre-orthogonalisation.
    All four methods cluster closely and track slightly above the
    raw baseline.}
  \label{fig:patho_auc_pca}
\end{figure}

\subsubsection{Aggregate Neuroscreening Performance}

Tables~\ref{tab:clinical_cv} and~\ref{tab:clinical_loso} report
aggregate neuroscreening MCC and AUC under 5-fold cross-validation
and LOSO respectively for all conditions.
Known results from the thesis baseline (unharmonised, site-wise
standardisation, and neuroCombat, 5-fold CV and LOSO) are included
where available; all other cells await completion of the extended
experiments.

\textbf{[TBF: narrative summary once all conditions are
complete, addressing: (i) whether harmonisation consistently improves
on the unharmonised baseline; (ii) additional aggregate benefit of
PCA pre-orthogonalisation; (iii) relative performance of CovBat
versus neuroCombat and site-wise standardisation; (iv) where the
DANN family sits relative to feature-level conditions; and (v)
concordance between LOSO and CV results.]}

\begin{table}[htbp]
  \centering
  \caption{Neuroscreening performance under 5-fold cross-validation
    (mean $\pm$ SD) for all conditions.
    $^\dagger$Signal-level models are evaluated under the LOSO
    protocol only (Section~\ref{subsubsec:clinical_eval}); 5-fold CV
    results are not applicable.}
  \label{tab:clinical_cv}
  \begin{tabular}{lcc}
    \toprule
    \textbf{Method} & \textbf{MCC} & \textbf{AUC} \\
    \midrule
    \multicolumn{3}{l}{\textit{GBE, without PCA pre-orthogonalisation}} \\
    \quad Unharmonised
      & $0.560 \pm 0.005$ & $0.863 \pm 0.001$ \\
    \quad Site-wise standardisation
      & $0.592 \pm 0.008$ & $0.878 \pm 0.003$ \\
    \quad neuroCombat
      & $0.590 \pm 0.011$ & $0.877 \pm 0.003$ \\
    \quad CovBat
      & \textbf{[TBF]} & \textbf{[TBF]} \\
    \midrule
    \multicolumn{3}{l}{\textit{GBE, with PCA pre-orthogonalisation}} \\
    \quad PCA only
      & \textbf{[TBF]} & \textbf{[TBF]} \\
    \quad PCA + site-wise stand.
      & \textbf{[TBF]} & \textbf{[TBF]} \\
    \quad PCA + neuroCombat
      & \textbf{[TBF]} & \textbf{[TBF]} \\
    \quad PCA + CovBat
      & \textbf{[TBF]} & \textbf{[TBF]} \\
    \midrule
    \multicolumn{3}{l}{\textit{Signal-level adversarial harmonisation$^\dagger$}} \\
    \quad MINet ($\lambda \equiv 0$)
      & \multicolumn{2}{c}{---} \\
    \quad MINet-MTL
      & \multicolumn{2}{c}{---} \\
    \quad MINet-DANN
      & \multicolumn{2}{c}{---} \\
    \quad MINet-DANN+tuning
      & \multicolumn{2}{c}{---} \\
    \bottomrule
  \end{tabular}
\end{table}

\begin{table}[htbp]
  \centering
  \caption{Neuroscreening performance under the Leave-One-Site-Out
    protocol (mean $\pm$ SD) for all conditions.}
  \label{tab:clinical_loso}
  \begin{tabular}{lcc}
    \toprule
    \textbf{Method} & \textbf{MCC} & \textbf{AUC} \\
    \midrule
    \multicolumn{3}{l}{\textit{GBE, without PCA pre-orthogonalisation}} \\
    \quad Unharmonised
      & $0.400 \pm 0.005$ & $0.737 \pm 0.006$ \\
    \quad Site-wise standardisation
      & $0.417 \pm 0.008$ & $0.752 \pm 0.007$ \\
    \quad neuroCombat
      & $0.427 \pm 0.011$ & $0.755 \pm 0.008$ \\
    \quad CovBat
      & \textbf{[TBF]} & \textbf{[TBF]} \\
    \midrule
    \multicolumn{3}{l}{\textit{GBE, with PCA pre-orthogonalisation}} \\
    \quad PCA only
      & \textbf{[TBF]} & \textbf{[TBF]} \\
    \quad PCA + site-wise stand.
      & \textbf{[TBF]} & \textbf{[TBF]} \\
    \quad PCA + neuroCombat
      & \textbf{[TBF]} & \textbf{[TBF]} \\
    \quad PCA + CovBat
      & \textbf{[TBF]} & \textbf{[TBF]} \\
    \midrule
    \multicolumn{3}{l}{\textit{Signal-level adversarial harmonisation}} \\
    \quad MINet ($\lambda \equiv 0$)
      & \textbf{[TBF]} & \textbf{[TBF]} \\
    \quad MINet-MTL
      & \textbf{[TBF]} & \textbf{[TBF]} \\
    \quad MINet-DANN
      & \textbf{[TBF]} & \textbf{[TBF]} \\
    \quad MINet-DANN+tuning
      & \textbf{[TBF]} & \textbf{[TBF]} \\
    \bottomrule
  \end{tabular}
\end{table}

% -----------------------------------------------------------------------
\section{Discussion}
\label{sec:discussion}

\subsection{The Nature and Magnitude of the Site Effect}
\label{subsec:disc_site_effect}

The site-classification experiment confirmed that the site effect is
both strong and broadly distributed across the 2\,850-dimensional
feature space of $\mathrm{ELM}_n$, with the CatBoost nonlinear
classifier achieving an overall MCC of $0.865 \pm 0.003$.
The range of per-site MCC values---from 0.954 (MOR) to 0.207
(CMD)---reflects the heterogeneity of the corpus: institutions with
distinctive demographic profiles (e.g.\ MOR and TOR, which serve
predominantly older patients) or unusual recording protocols
(e.g.\ KATMOJPRZ, with sessions up to 8\,322 seconds) are most
reliably identified.
Institutions at the lower end of the MCC range (CMD, KIEG, CHE,
SL2) are likely to have feature distributions that overlap
substantially with those of neighbouring sites, either because of
similar patient demographics or because their recording protocols
are less distinctive.

The SHAP analysis and ablation study converge on a consistent
interpretation: inter-electrode coherence features are the primary
carriers of site-specific information.
The elevated importance of O1--O2 coherence is consistent with
the hypothesis that differences in eye-open versus eye-closed
protocols across institutions modulate occipital alpha coherence,
which is highly sensitive to visual state.
The prominence of beta-band coherence in temporal regions is
consistent with variability in muscle artefact levels arising from
differences in patient preparation, scalp cleaning, and electrode
impedance management across sites, since the beta band overlaps
with the electromyographic spectrum.
Delta-band power and frontal Fz power, identified as important
spectral components, may reflect differences in patient alertness
and vigilance state at recording time.

\subsection{The Harmonisation Paradox and the Two-Classifier
  Interpretation}
\label{subsec:disc_paradox}

The most striking finding of this study is the \emph{harmonisation
paradox}: all marginal correction methods successfully reduced
pairwise Cohen's~$d$ inter-site differences and drove the logistic
regression linear probe to near-chance site-classification MCC, yet
simultaneously increased the CatBoost nonlinear probe MCC towards a
near-ceiling value.

The two-classifier framework resolves this apparent contradiction.
Marginal correction adjusts the mean and variance of each feature
independently, thereby equalising marginal distributions across
sites and removing the linear component of the site effect.
However, the site effect is encoded not only in marginal
distributions but also in the \emph{covariance structure} between
features---a form of inter-site variability that feature-by-feature
corrections leave entirely
intact~\cite{hu2023image}.
By equalising means and variances, these methods remove the dominant
sources of within-site variance, causing the subtler but highly
structured inter-feature covariance patterns that encode site
identity to become more salient to a nonlinear classifier.
This argument is directly analogous to that advanced by
Mellot~et~al.~\cite{mellot2023meeg} in the context of Riemannian
geometry-based EEG alignment.

The PCA pre-orthogonalisation experiments provide direct mechanistic
evidence: by projecting the feature space onto mutually uncorrelated
principal components, PCA eliminates inter-feature covariance in the
original feature basis.
Marginal correction applied to this PC-transformed space therefore
cannot amplify covariances in the original basis, and the gap between
LR-MCC and CatBoost-MCC narrows substantially.

This finding has a direct methodological implication: univariate
metrics such as Cohen's~$d$ and the linear probe MCC are inadequate
as sole indicators of harmonisation quality in high-dimensional
feature spaces.
A complete evaluation framework must include a nonlinear
site-discriminability measure---such as the CatBoost MCC reported
here---as a necessary complement to feature-level effect-size
summaries.

\subsection{Signal-Level Adversarial Harmonisation}
\label{subsec:disc_dann}

The four DANN variants evaluated in this study are designed to test
whether end-to-end representation learning can overcome the
fundamental limitation of marginal correction.
The theoretical prediction is clear: MINet-MTL, which trains a domain
classifier without gradient reversal, should actively exploit site
identity and is therefore expected to produce a larger CatBoost-MCC
on $\mathbf{h}$ than MINet, representing an empirical upper bound on
the site information an unconstrained multi-task network encodes in
the recording-level representation.
MINet-DANN, which trains with the GRL, should suppress site
information at the representation level; the key empirical question
is whether this suppression extends to the \emph{nonlinear}
covariance component, as measured by the CatBoost probe on
$\mathbf{h}$.
MINet-DANN+tuning addresses a practical limitation of adversarial
training: the learned representation may not immediately yield
well-calibrated diagnostic predictions for a new site, and
recalibrating the pathology classifier $G_y$ on a frozen,
site-agnostic backbone recovers diagnostic utility without
modifying the harmonised representation.

\textbf{[TBF: complete the discussion of DANN results once all
experimental conditions are finalised, addressing: (i) whether
MINet-DANN reduces the LR--CatBoost gap on $\mathbf{h}$ relative to
all feature-level conditions; (ii) whether MINet-DANN+tuning recovers
or exceeds the neuroscreening performance of the best feature-level
condition under LOSO; and (iii) the theoretical implications of the
observed LR--CatBoost gap pattern across all four DANN variants.]}

\subsection{Impact on Clinical Neuroscreening}
\label{subsec:disc_clinical}

\textbf{[TBF: complete once all conditions are finalised, addressing
concordance between CV and LOSO results, the relative benefit of PCA
pre-orthogonalisation on clinical performance, the position of the
DANN family relative to feature-level methods, and the practical
significance of the observed AUC and MCC improvements for clinical
deployment.]}

\subsection{Limitations}
\label{subsec:limitations}

Several limitations of the present study should be noted.
First, although patient age and sex were included as biological
covariates in neuroCombat and CovBat, other clinically relevant
variables---such as the reason for referral, medication status, or
cognitive state at recording time---were not available for the full
corpus and could not be controlled for.
Second, all marginal correction methods adjust only the first two
moments of each feature's distribution and cannot correct
higher-order distributional differences.
Third, the LOSO calibration procedure (using $n_{\mathrm{cal}} = 30$
normal recordings from the held-out site) introduces a practical
dependency on the availability of a small representative sample from
any new site.
For feature-level methods this sample is used solely to fit
harmonisation parameters; for signal-level DANN variants it is
included in the source-pool training data, introducing a minor
asymmetry relative to the strict LOSO setting.
Given that $n_{\mathrm{cal}} = 30$ constitutes fewer than 0.06\% of
the total training corpus, the expected bias on reported performance
estimates is negligible, but this should be borne in mind when
comparing feature-level and signal-level results directly.
Fourth, PCA pre-orthogonalisation discards the covariance structure
of the original feature space; while this attenuates the
harmonisation paradox, it also removes inter-feature relationships
that may carry clinically relevant information, and the long-term
impact on downstream task performance warrants further investigation.
Fifth, the CatBoost probe applied to the DANN representations
$\mathbf{h}$ evaluates nonlinear site discriminability in a
288-dimensional latent space rather than the original 2\,850-dimensional
feature space; while this is internally consistent within the
two-classifier framework, direct numerical comparison between the
feature-level and signal-level CatBoost-MCC values should be
interpreted with caution.

\subsection{Future Directions}
\label{subsec:future}

The harmonisation paradox and the two-classifier framework identified
in this study motivate several directions for future work.
Most immediately, the completed DANN results will determine whether
adversarial training can suppress nonlinear covariance-based site
structure, which neither marginal correction nor PCA
pre-orthogonalisation can fully address.
Beyond harmonisation methodology, the SHAP analysis suggests that
feature engineering strategies that reduce coherence-feature
sensitivity to site-specific procedural differences could reduce the
effective site effect without explicit harmonisation.
Finally, normative modelling frameworks~\cite{dong2025normative},
which express individual measurements as deviations from an
age-adjusted population norm, constitute a complementary approach
that may interact favourably with the harmonisation methods studied
here and warrant systematic evaluation in the $\mathrm{ELM}_s$
context.

% -----------------------------------------------------------------------
\section{Conclusion}
\label{sec:conclusion}

This study provides a comprehensive characterisation of the site
effect in a large-scale clinical EEG corpus ($\mathrm{ELM}_s$,
54\,779 recordings, 30 hospitals) and a systematic evaluation of
feature-level and signal-level harmonisation strategies under a
unified two-classifier evaluation framework.

The CatBoost nonlinear classifier achieved a site-identification MCC
of $0.865 \pm 0.003$, confirming that site identity is strongly
encoded in EEG feature representations.
SHAP explainability analysis and ablation experiments identified
inter-electrode coherence---particularly O1--O2 and beta-band
temporal coherence---as the dominant carrier of site-specific
information, implicating differences in eye-state protocols and
muscle artefact management across institutions.

A central finding is the \emph{harmonisation paradox}: all marginal
correction methods successfully removed the linear component of the
site effect (Cohen's~$d$ reduced; logistic regression MCC near
chance) yet simultaneously amplified the nonlinear CatBoost MCC
to near-ceiling values, demonstrating that feature-by-feature
correction reorganises rather than eliminates the multivariate
covariance structure of the site effect.
PCA pre-orthogonalisation attenuates this paradox and yields more
uniform and stable cross-site neuroscreening performance.
These findings demonstrate that a complete harmonisation evaluation
framework must include a nonlinear site-discriminability measure
alongside univariate metrics.

\textbf{[TBF: summarise DANN and full neuroscreening results once
all experimental conditions are complete.]}

These results collectively underscore that the site effect is a
complex, high-dimensional challenge for which standard marginal
correction methods provide only a partial solution.
Developing truly generalisable neuroscreening models will require
harmonisation strategies that explicitly account for multivariate
inter-site structure---whether through covariance-aware statistical
corrections, adversarial representation learning, or feature
engineering guided by explainability analysis.

% -----------------------------------------------------------------------
\section*{Acknowledgements}
The authors are grateful to Elmiko Biosignals sp.~z~o.o.\ for
providing access to the ELM19 database.

\section*{CRediT Author Contribution Statement}
K.C.: Methodology, Software, Formal Analysis, Investigation,
Visualisation, Writing -- Original Draft.
J.P.: Methodology, Writing -- Review \& Editing.
J.\.Z.: Conceptualisation, Supervision, Resources, Writing --
Review \& Editing.

\section*{Declaration of Competing Interests}
The authors declare that they have no known competing financial
interests or personal relationships that could have appeared to
influence the work reported in this paper.

\section*{Data Availability}
The ELM19 database is available from Elmiko Biosignals sp.~z~o.o.\
upon reasonable request to
\href{mailto:EEG.database@elmiko.pl}{EEG.database@elmiko.pl}.

% -----------------------------------------------------------------------
\appendix
\section{Patient Demographics by Institution}
\label{app:demographics}

\begin{table}[htbp]
  \centering
  \caption{Summary statistics of patient age (years) by institution
    for the $\mathrm{ELM}_n$ normal recording subset
    (representative selection).}
  \label{tab:age_stats}
  \begin{tabular}{lrrrrr}
    \toprule
    \textbf{Institution} & \textbf{Mean} & \textbf{Median} &
    \textbf{SD} & \textbf{Min} & \textbf{Max} \\
    \midrule
    PIO    & 20.3 & 17.6 & 8.3  & 16.0 & 68.3 \\
    SL2    & 28.4 & 19.5 & 15.1 & 16.0 & 66.3 \\
    TER\_L & 26.3 & 17.8 & 15.0 & 16.0 & 69.8 \\
    MKW    & 26.1 & 21.7 & 10.8 & 16.0 & 69.7 \\
    GAK    & 41.3 & 40.9 & 13.8 & 16.0 & 70.0 \\
    KAL    & 46.2 & 48.2 & 15.6 & 17.3 & 69.9 \\
    OST    & 48.1 & 49.2 & 13.7 & 18.4 & 69.9 \\
    MOR    & 48.6 & 52.1 & 15.4 & 18.7 & 69.4 \\
    TOR    & 47.6 & 50.8 & 16.5 & 16.0 & 69.9 \\
    \multicolumn{6}{c}{$\cdots$} \\
    \bottomrule
  \end{tabular}
\end{table}

% -----------------------------------------------------------------------
\bibliographystyle{elsarticle-num}
\bibliography{bibliography}

\end{document}